\section{After The Needle}

The first chapter ended with a narrow permission. After the collapse of
truth, the grammar could not recover by announcing a universe of ready-made
objects. It could only sanction one local identity, a needle strong enough
to let a reading return to itself and weak enough not to smuggle in a whole
metaphysics. That needle is the hinge on which the present chapter turns.
With it, the grammar can stabilize truth long enough to ask what a finite
measurement must preserve when it is varied.

The answer is not yet a field, a force, or a law of nature. The answer is a
gauge: a finite rule for comparing neighboring readings while forbidding
the comparison from pretending to see beyond its own horizon. A gauge path
carries a sequence of permitted distinctions. An action assigns a cost to
that sequence. A variation tests whether an admissible slip changes the
reading. The whole construction remains finite. It may refine, but each
refinement must leave a record that the grammar can read.

This finite setting changes the meaning of stationarity. A stationary
reading is not a picture of something frozen. It is a path whose permitted
local slips leave no first residual to spend. The grammar reads the path,
reads the allowed variations, and finds that the residue vanishes. Because
the residue is the only authorized witness of imbalance, its vanishing is
the grammatical form of stationarity.

The chapter follows that single obligation through six renderings. First,
finite gauge paths and actions give stationarity its local reader. Second,
the vanishing first variation and the discrete Euler--Lagrange equations
become two ways to read one constraint. Third, a closed balanced cubic
spline shows that finite anchors can force the remaining residue to zero.
Fourth, real partitions and squeezed residuals explain how finite
truncations approach completion without pretending that the continuum was
given in advance. Fifth, a boundary anchor kills the nullspace that would
otherwise let unread activity pass as zero. Finally, the second variation is
identified as the unique positive invariant carried through the whole
chain.

The title of the chapter is therefore literal. A finite gauge is enough to
read one invariant, but only after the grammar has forbidden every cheaper
escape: unearned continuity, free boundary drift, invisible null modes, and
residual imbalance. What remains is not a decorative scalar. It is the
single quantity the grammar is forced to carry.
